% !TEX program = xelatex

%page setup
\documentclass[svgnames,x11names,b5paper]{book}
\usepackage[b5paper, portrait, margin=1.5cm]{geometry}

%font shit
\usepackage{fontspec}
\setmainfont{Libertinus Serif}

%for typing in different scripts
\usepackage{polyglossia}
\renewcommand{\rm}{\normalfont}
%\newfontfamily\russianfont{MonomakhUnicode}
\setdefaultlanguage{english}
\setotherlanguages{armenian,sanskrit,greek,russian}

%\setlength{\parindent}{10pt}
%random stuff which handles tables
\usepackage{array}
\usepackage{multirow}
\usepackage{multicol}
\usepackage{float}

%asmsymbol
\usepackage{amssymb}

%lists
\usepackage{enumitem}
\setlist{noitemsep}

%pagecolour
\usepackage{xcolor}
%\pagecolor{AntiqueWhite}

%tikz
\usepackage{tikz}
\usepackage{tikz-qtree}

%footnote formatting
\usepackage[bottom,hang,multiple]{footmisc}
\renewcommand\footnoterule{}
\setlength{\footnotemargin}{0.75em}

%footnotes in tables
%\newcommand{\fnark}{\addtocounter{footnote}{1}\footnotemark[\thefootnote]}
%\newcommand{\fnext}[1]{\footnotetext{#1}}

% baarux
\usepackage[mcolblock,nostandards]{baabbrevs}
\usepackage[canonical,nochapterlabels]{baarux-0.9.9}
\baaruset{junctureoutersep=0.4ex, junctureinnersep=0.2ex}
\baaruset{hangindent=1.5em} % adjusting indenting on multiline glosses
\renewcommand{\glossarypreamble}{\thispagestyle{empty}} % for removing page number from abbreviations glossary
% so to make an abbreviation work in the gloss, you need to make a command \baabbrev{abbreviation}{explanation to put in the glossary}

\baabbrev[\FIRST]{1}{first person}
\baabbrev[\SECOND]{2}{second person}
\baabbrev[\THIRD]{3}{third person}
\baabbrev[\OBJ]{o}{object}
\baabbrev[\SUBJ]{s}{subject}

%table of contents
\usepackage{titletoc}
\usepackage{titlesec}
\setcounter{tocdepth}{4}

%random formatting mumbo-jumbo
%\usepackage{hyperref}
\usepackage{ragged2e}

% subsubsubsections
\setcounter{secnumdepth}{4}
\titleformat{\paragraph}
{\normalfont\normalsize\bfseries}{\theparagraph}{1em}{}
\titlespacing*{\paragraph}{0pt}{3.25ex plus 1ex minus .2ex}{1.5ex plus .2ex}

\usepackage{booktabs}

%macro that allows you to change your lang's name easily; when in prose type \langname{}
\newcommand{\langname}{Carite}

% for making dictionary entries:
% in the first {} put the lexeme
% in the second put the pronunciation
% in the third put the part of speech
% in the fourth put stuff like gender, transitivity, whatever
% in the fifth put principal parts or genitives of nouns etc
% in the sixth put the meanings
% in the seventh put etymological information
% to cite a source within a dictentry, put \fnark{} where you want the footnote to be and then put \fnext{} after the entry, and if you have more than one footnote in the entry you need to put \addtocounter{footnote}{-n} for every n+1 footnotes you add

\newcommand{\dictentry}[7]{

\begin{minipage}{\textwidth}
\textbf{#1} \quad{} {/}{#2}{/} \quad{} \textsc{#3}. \quad{} \textsc{#4} ~\textit{#5} \\\vspace{4pt}
\textsc{senses} #6 \\\vspace{-2pt}
\textsc{etym} ~~ #7
\end{minipage}\\\vspace{8pt}

}

\begin{document}

% !TEX program = xelatex

{\pagestyle{empty}

\centering

\null\vfill

\fontsize{20}{24}\selectfont

\textsc{A title or something} \vspace{12pt}

\hline

\vspace{12pt}

\textsc{Jasper}\vspace{240pt}

\vfill

\clearpage}

\newpage

\tableofcontents
\clearpage

\section{Introduction}

\section{Sound Changes}

%% for sc in lt_soundchanges
    %% if not 'donotdisplay' in sc
        %% if 'newchapter' in sc
            \subsection{\VAR{sc.newchapter}}
        
        %% endif
        \VAR{sc.prose}

        %% if 'examples' in sc
            \begin{table}[H]
                \begin{center}
                \begin{tabular}{llcllllll}
            %% for example in sc.examples
                \VAR{example}\\
            %% endfor
                \end{tabular}
                \end{center}
                \end{table}
        
        %% endif
        %% if 'prose_after' in sc
            \VAR{sc.prose_after}

        %% endif
    %% endif
%% endfor

\begin{table}[H]
    \begin{center}
    \begin{tabular}{lllllllll}
    \textsc{PIE} & *sem- & → & \textsc{proto-Hellenic} & *hen-s & → & \textsc{Attic} & εἵς \textit{heís} & `one'
    \end{tabular}
    \end{center}
    \end{table}

\section{Phonology}

\section{Morphology}

% you can get things to span n columns with \multicolumn{n}{c}{blabla}
\begin{table}[H]
    \begin{center}
    \begin{tabular}{r|ccc}
    \toprule
    & \textsc{sg} & \textsc{du} & \textsc{pl} \\
    \midrule
    \textsc{nom} & λόγος & \multirow{2}{*}{λόγω} & λόγοι\\
    \textsc{acc} & λόγον & & λόγους\\
    \textsc{gen} & λόγου & \multirow{2}{*}{λόγοιν} & λόγων\\
    \textsc{dat} & λόγωι & & λόγοις\\
    \bottomrule
    \end{tabular}
    \end{center}
    \end{table}

\section{Lexicon}

~~~~
% ^ these tildes allow for indentation and prevent the first dictionary entry from being too far left

%% for entry in lt_dictionary
\dictentry{\VAR{entry.spelling}}{\VAR{entry.ipa}}{\VAR{entry.wordclass}}{q-stem}{ataʀum}{\VAR{entry.meaning}}{\VAR{entry.etymo}}
%% endfor

prose prose prose\\

\end{document}